\section{Introduction}
Approximate machine unlearning is intended to remove selected training influence without
the cost of complete retraining. The relevant counterfactual is a model trained on the
retained data, but stochastic optimization produces a distribution of such models. This
observation changes the evaluation problem. A single retrained network is one draw from a
random training procedure, not a unique parameter-space target. Differences caused by
initialization, minibatch order, optimization, and augmentation can be comparable to or
larger than those caused by deleting a small set of observations.

Many empirical evaluations emphasize retained-set utility, decreased performance on a
forget set, or membership-inference behavior. Each measurement can be useful, but none is
identical to the counterfactual question: would the observable and latent behavior of the
post-unlearning model be statistically consistent with models trained without the deleted
records? A model can preserve test accuracy while changing its confidence distribution; an
attack can approach chance while the model remains behaviorally unlike retraining; and an
IID-deleted example can remain correctly classified even when it was never used for
training. These possibilities motivate a multidimensional target rather than a single
forgetting score.

This study separates two questions that are often combined. First, is the deletion request
itself empirically identifiable under the selected model, data, and metrics? Second,
conditional on such sensitivity, does an approximate method reproduce the distribution of
exactly retrained models? If the original model cannot be distinguished from the reference
population, approximate/exact proximity is compatible both with genuine removal and with a
procedure that leaves the original function essentially unchanged. In that setting the
assay cannot support a positive or negative method claim.

The present manuscript reports a bounded retrospective pilot, P001, and the infrastructure
used to preserve it. P001 uses one reduced Fashion-MNIST split, one frozen IID deletion
manifest, one original seed, and 15 exact-retraining seeds. Its purpose is methodological:
to validate exclusion, lineage, immutable evidence, reference-ensemble estimation, and
paper generation. The analysis then asks whether P001 was an informative deletion assay.
It makes no novelty or successful-unlearning claim.

The distinction between an experimental target and an evaluation convenience is important in
this setting. Let $D$ be a training set, $F\subset D$ a requested deletion set, $A$ a
stochastic training procedure, and $U$ an approximate update procedure. The usual objects are
an original model $M_0=A(D)$, an updated model $M_U=U(M_0,F)$, and one retrained model
$M_R=A(D\setminus F)$. The latter is often used as though it were a fixed oracle. It is not.
For a fixed retained dataset, $A(D\setminus F)$ induces a distribution over functions,
probabilities, calibration states, and representations. A scientifically cautious comparison
therefore asks whether $M_U$ is compatible with that distribution on prespecified probes. It
does not ask whether $M_U$ happens to resemble one convenient retraining draw.

This perspective also clarifies why a favorable conventional outcome can be ambiguous. A
reduction in forget-set accuracy may reflect meaningful removal, ordinary class confusion, or
collateral degradation. Stable retained accuracy may coexist with changed confidence or local
behavior. A weak membership attack may be weak because of its threat model rather than because
the model is counterfactually faithful. None of these observations is useless; each answers a
different question. The manuscript therefore treats their agreement and disagreement as an
object of measurement, while declining to collapse them into one unsupported success score.

The project is intentionally staged. Tier-0 and Tier-1 work establishes deterministic data
handling, artifact contracts, and deletion invariants. The current P001 release is a bridge
between those engineering checks and a future publication-scale study. It is designed to make
the transition auditable: what was measured, what was merely designed, what failed, and what
remains unmeasured are all separated in the repository and in the text. That separation is less
dramatic than a method leaderboard, but it is necessary if later empirical claims are to be
interpretable under skeptical review.

The contributions of this release are therefore deliberately narrow:
\begin{enumerate}
\item an operational counterfactual evaluation model in which exact retraining is an
ensemble rather than one arbitrarily selected network;
\item an audit pipeline that cryptographically binds dataset, split, deletion, training,
prediction, metric, figure, and manuscript artifacts;
\item a reproducible P001 case study demonstrating seed-count diagnostics and the need to
test deletion identifiability before interpreting method proximity; and
\item a claim boundary and prospective C001 gate that prevent exploratory pilot outcomes
from being presented as confirmatory evidence.
\end{enumerate}

The remainder of the paper first positions this scope relative to unlearning, auditing, and
evaluation work. It then specifies the counterfactual target, evidence lineage, measurement
layers, and statistical boundaries. The results section reports only the completed pilot and
its negative assay-identifiability finding. Later sections describe the implications, failure
records, limitations, and the conditions that must be met before a publication-grade
confirmatory study is allowed to start.
